% Part: normal-modal-logic
% Chapter: frame-definability
% Section: definability

\documentclass[../../../include/open-logic-section]{subfiles}

\begin{document}

\olfileid{nml}{frd}{def}

\olsection{చట్రాల నిర్వచనీయత}

\olref[acc]{thm:soundschemas}లోని విలోమ సూచనలు
నమూనాల విషయంలో విఫలమైనా, ``నమూనా'' స్థానంలో
``చట్రం''ను పెడితే నిలుస్తాయి: \olref[acc]{thm:soundschemas}లో
పరిశీలించిన ధర్మాల విషయంలో, ఒక \emph{చట్రంలో}
\tetoken{సూత్రం}{!!a{formula}} చెల్లుబాటు అయితే ఆ చట్రం
ప్రాప్యత సంబంధానికి అనురూపమైన ధర్మం ఉంటుందన్నది
\emph{నిజం}. కాబట్టి పరిశీలించిన \tetoken{సూత్రాలు}{!!{formula}s}
అనురూప ధర్మం గల చట్రాల వర్గాలను \emph{నిర్వచిస్తాయి}.

\begin{defn}
  $\mClass{F}$ చట్రాల వర్గమైతే, ఖచ్చితంగా
  $\mModel{F} \in \mClass{F}$ అయ్యే చట్రాలన్నింటికీ,
  వాటికే, $\mModel{F} \Entails !A$ ఉన్నప్పుడే,
  అప్పుడే $!A$ అనేది~$\mClass{F}$ను
  \emph{నిర్వచిస్తుంది} అంటాం.
\end{defn}

ఇప్పుడు చట్రాల నిర్వచనీయతకు సంబంధించిన పూర్తి
ఫలితాలను స్థాపిద్దాం.

\begin{thm}\ollabel{thm:fullCorrespondence}
\olref[acc]{tab:five} కుడివైపు ఉన్న \tetoken{సూత్రం}{!!{formula}}
చట్రం~$\mModel{F}$లో చెల్లుబాటు అయితే, ఎడమవైపు
ఉన్న ధర్మం $\mModel{F}$కు ఉంటుంది.
\end{thm}

\begin{proof}
  \begin{enumerate}
  \item \Ax{D} అనేది $\mModel{F} = \tuple{W, R}$లో
    చెల్లుబాటు అవుతుందని, అంటే
    $\mModel{F} \Entails \Box p \lif \Diamond p$ అని
    అనుకుందాం. $\mModel{F}$పై ఆధారపడే
    $\mModel{M} = \tuple{W, R, V}$ నమూనాను, $w \in W$ను
    తీసుకుందాం. $Rwv$ అయ్యే ఒక $v$ ఉందని చూపాలి.
    అలా లేదనుకుందాం: అప్పుడు $p$తో సహా ఏ~$!A$కైనా,
    $\mSat{M}{\Box !A}[w]$ మరియు
    $\mSat/{M}{\Diamond !A}[w]$.
    \sourcecorrection{OLTENMLFRDDEF-001}{మూలంలో మొదటి
      సంతృప్తి సంకేతానికి లోకం $[w]$ లేదు. ప్రాప్య లోకాలు
      లేనప్పుడు $\Box !A$ అనేది ఈ $w$ వద్ద సత్యం;
      మొత్తం నమూనాలో తప్పనిసరిగా సత్యం కాదు. పక్కనున్న
      డైమండ్ సంకేతంతో సమానంగా లోక సూచిని చేర్చాం.}
    అయితే $\mSat/{M}{\Box p \lif
      \Diamond p}[w]$; ఇది $\mModel{F}
    \Entails \Box p \lif \Diamond p$ అన్న ఊహకు వైరుధ్యం.
  \item \Ax{T} అనేది $\mModel{F}$లో చెల్లుబాటు
    అవుతుందని, అంటే $\mModel{F} \Entails \Box
    p \lif p$ అని అనుకుందాం. $w \in W$ అనే ఏకపక్ష
    లోకాన్ని తీసుకుందాం; $Rww$ చూపాలి. $Rwu$
    అయినప్పుడే, అప్పుడే $u \in V(p)$ అని కేటాయిద్దాం
    ($q$ అనేది $p$ కాకుండా వేరే చరమైతే, $V(q)$
    ఏదైనా కావచ్చు; ఉదాహరణకు $V(q) = \emptyset)$.
    $\mModel{M} = \tuple{W, R, V}$ అనుకుందాం.
    నిర్మాణం వల్ల $Rwu$ అయ్యే అన్ని $u$లకు
    $\mSat{M}{p}[u]$; అందువల్ల
    $\mSat{M}{\Box p}[w]$. ఊహ ప్రకారం $\Box p \lif p$
    అనేది~$w$ వద్ద సత్యం కాబట్టి $\mSat{M}{p}[w]$.
    కానీ $V$ నిర్వచనం ప్రకారం ఇది $Rww$ అయినప్పుడే
    సాధ్యం.
  \item వ్యతిరేక ప్రతిజ్ఞను నిరూపిద్దాం: $\mModel{F}$
    సౌష్ఠవం కాదని అనుకొని, \Ax{B}, అంటే
    $p \lif \Box\Diamond p$, అనేది
    $\mModel{F}= \tuple{W, R}$లో చెల్లుబాటు కాదని
    చూపుదాం. $\mModel{F}$ సౌష్ఠవం కాకపోతే
    $Ruv$ ఉండి $Rvu$ లేని $u$, $v \in W$
    ఉంటాయి. $Rvw$ కానప్పుడే, అప్పుడే $w \in V(p)$
    అయ్యేలా $V$ను నిర్వచిద్దాం (ఇతర చరాలపై
    $V$ ఏదైనా కావచ్చు). $\mModel{M} =\tuple{W, R,
      V}$ అనుకుందాం. $V$ నిర్వచనం వల్ల $Rvw$
    కాని అన్ని $w$ల వద్ద $\mSat{M}{p}[w]$;
    ముఖ్యంగా $Rvu$ కాదు కాబట్టి $\mSat{M}{p}[u]$.
    అంతేకాక, $Rvw$ అప్పుడే, అప్పుడే
    $w \notin V(p)$ కాబట్టి $Rvw$ మరియు
    $\mSat{M}{p}[w]$ రెండూ అయ్యే $w$ ఏదీ లేదు;
    కనుక $\mSat/{M}{\Diamond p}[v]$. $Ruv$
    కూడా ఉంది కాబట్టి
    $\mSat/{M}{\Box\Diamond p}[u]$.
    అందువల్ల $\mSat/{M}{p \lif
      \Box\Diamond p}[u]$; కాబట్టి $\Ax{B}$
    అనేది~$\mModel{F}$లో చెల్లుబాటు కాదు.
  \item \Ax{4} అనేది $\mModel{F} = \tuple{W,R}$లో
    చెల్లుబాటు అవుతుందని, అంటే
    $\mModel{F} \Entails \Box p \lif \Box\Box p$ అని
    అనుకుందాం. $Ruv$, $Rvw$ అయ్యే ఏకపక్ష
    లోకాలు $u$, $v$, $w \in W$ను తీసుకుందాం;
    $Ruw$ చూపాలి. $Ruz$ అయినప్పుడే, అప్పుడే
    $z \in V(p)$ అయ్యేలా $V$ను నిర్వచిద్దాం
    (మిగిలిన చరాలపై $V$ ఏదైనా కావచ్చు).
    $\mModel{M} =\tuple{W, R, V}$ అనుకుందాం.
    $V$ నిర్వచనం వల్ల $Ruz$ అయ్యే అన్ని $z$లకు
    $\mSat{M}{p}[z]$; కనుక $\mSat{M}{\Box p}[u]$.
    అయితే ఊహ ప్రకారం \Ax{4}, అంటే $\Box p
    \lif \Box \Box p$, అనేది $u$ వద్ద సత్యం;
    అందువల్ల $\mSat{M}{\Box \Box
      p}[u]$. $Ruv$, $Rvw$ కాబట్టి
    $\mSat{M}{p}[w]$; $V$ నిర్వచనం ప్రకారం
    ఇది $Ruw$ అయినప్పుడే సాధ్యం. అదే కావలసింది.
  \item వ్యతిరేక ప్రతిజ్ఞ పద్ధతిలో, చట్రం
    $\mModel{F} = \tuple{W, R}$ యూక్లిడియన్ కాదని
    అనుకొని, అది \Ax{5}ను అసత్యం చేస్తుందని,
    అంటే $\mModel{F} \Entails/ \Diamond p \lif
      \Box\Diamond p$ అని చూపుదాం. $Rwu$,
    $Rwv$ ఉండి $Ruv$ లేని లోకాలు $u$, $v$,
    $w \in W$ ఉన్నాయనుకుందాం. అన్ని లోకాలు $z$కు,
    $Ruz$ \emph{కాని}ప్పుడే, అప్పుడే
    $z \in V(p)$ అయ్యేలా $V$ను నిర్వచిద్దాం.
    $\mModel{M} =\tuple{W, R, V}$ అనుకుందాం.
    అప్పుడు ఊహ ప్రకారం $\mSat{M}{p}[v]$;
    $Rwv$ కూడా కాబట్టి $\mSat{M}{\Diamond p}[w]$.
    కానీ $Ruy$, $\mSat{M}{p}[y]$ రెండూ అయ్యే లోకం
    $y$ లేదు; కాబట్టి $\mSat/{M}{\Diamond
      p}[u]$. $Rwu$ కాబట్టి
    $\mSat/{M}{\Box\Diamond
      p}[w]$. అందువల్ల \Ax{5}, అంటే
    $\Diamond p \lif \Box\Diamond p$, అనేది~$w$
    వద్ద విఫలమవుతుంది.
  \end{enumerate}
\end{proof}

\Ax{D} నిరూపణకీ ఇతర సందర్భాల నిరూపణలకీ
ఒక తేడా గమనించవచ్చు: అందులో మూల్యనిర్ణయం~$V$
గురించి ప్రస్తావించలేదు. నిజానికి,
$\mSat{M}{\Ax{D}}$ అయితే $\mModel{M}$
సీరియల్ అని నిరూపించాం. అందువల్ల \Ax{D}
సీరియల్ \emph{నమూనాల} వర్గాన్నే నిర్వచిస్తుంది;
చట్రాల వర్గాన్ని మాత్రమే కాదు.

\begin{cor}\ollabel{prop:D-serial}
  \Ax{D} సత్యమైన ప్రతి నమూనా సీరియల్.
\end{cor}

\begin{cor}
\olref[acc]{tab:five} కుడివైపు ఉన్న ప్రతి
\tetoken{సూత్రం}{!!{formula}} ఎడమవైపు ధర్మం గల
చట్రాల వర్గాన్ని నిర్వచిస్తుంది.
\end{cor}

\begin{proof}
  \olref[acc]{thm:soundschemas}లో, ఒక నమూనాకు
  ఎడమవైపు ధర్మం ఉంటే, కుడివైపు
  \tetoken{సూత్రం}{!!{formula}} అందులో సత్యమని
  నిరూపించాం. కాబట్టి చట్రం~$\mModel{F}$కు
  ఎడమవైపు ధర్మం ఉంటే, కుడివైపు
  \tetoken{సూత్రం}{!!{formula}} $\mModel{F}$లో
  చెల్లుబాటు అవుతుంది. \olref{thm:fullCorrespondence}లో
  విలోమ సూచనలను నిరూపించాం: కుడివైపు ఉన్న
  \tetoken{సూత్రం}{!!a{formula}} చట్రం~$\mModel{F}$లో
  చెల్లుబాటు అయితే, ఎడమవైపు ధర్మం
  $\mModel{F}$కు ఉంటుంది.
\end{proof}

\begin{prob}
\olref[nml][frd][acc]{tab:anotherfive} కుడివైపు ఉన్న
\tetoken{సూత్రం}{!!{formula}} చట్రం~$\mModel{F}$లో
చెల్లుబాటు అయితే, ఎడమవైపు ధర్మం $\mModel{F}$కు
ఉంటుందని చూపండి. దీనికోసం ఎడమవైపు ధర్మం
\emph{లేని} చట్రాన్ని తీసుకొని, కుడివైపు
\tetoken{సూత్రం}{!!{formula}} ఏదో లోకంలో అసత్యమయ్యేలా
తగిన~$V$ను నిర్వచించండి.
\end{prob}

\olref{thm:fullCorrespondence}తో ధర్మాలను కలిపి
పరిగణించవచ్చని కూడా తెలుస్తుంది: ఉదాహరణకు,
\Ax{B}, \Ax{4} రెండూ~$\mModel{F}$లో
చెల్లుబాటు అయితే ఆ చట్రం సౌష్ఠవమూ,
సంక్రామకమూ అవుతుంది. కొన్ని చట్ర ధర్మాలను
కలిపి కలిగిన అన్ని చట్రాల్లో చెల్లుబాటు అయ్యే
\tetoken{సూత్రాల}{!!{formula}s} సమితిగా అనేక ముఖ్యమైన
మోడల్ తర్కాలను లక్షణీకరిస్తారు. కనుక వాటిని,
అనురూప నిర్వచక \tetoken{సూత్రాలు}{!!{formula}s}
చెల్లుబాటు అయ్యే అన్ని చట్రాల్లో చెల్లుబాటు అయ్యే
\tetoken{సూత్రాల}{!!{formula}s} సమితిగా కూడా లక్షణీకరించవచ్చు.
ఉదాహరణకు, సాంప్రదాయిక వ్యవస్థ \Log{S4} అనేది
అన్ని స్వావర్తన, సంక్రామక చట్రాల్లో చెల్లుబాటు
అయ్యే \tetoken{సూత్రాల}{!!{formula}s} సమితి;
అంటే \Ax{T},~\Ax{4} రెండూ చెల్లుబాటు అయ్యే
అన్ని చట్రాల్లో చెల్లుబాటు అయ్యేవి.
\Log{S5} అనేది అన్ని స్వావర్తన, సౌష్ఠవ,
యూక్లిడియన్ చట్రాల్లో చెల్లుబాటు అయ్యే అన్ని
సూత్రాల సమితి; అంటే \Ax{T}, \Ax{B}, \Ax{5}
మూడూ చెల్లుబాటు అయ్యే అన్ని చట్రాల్లో
చెల్లుబాటు అయ్యే సూత్రాల సమితి.

$R$ ధర్మాల మధ్య తార్కిక సంబంధాలకు సాధారణంగా
వాటి అనురూప నిర్వచక \tetoken{సూత్రాల}{!!{formula}s}
మధ్య సంబంధాలు అనుగుణంగా ఉంటాయి. ఉదాహరణకు,
ప్రతి స్వావర్తన సంబంధం సీరియల్ కూడా; అందువల్ల
ఒక చట్రంలో \Ax{T} చెల్లుబాటు అయితే~\Ax{D}
కూడా చెల్లుబాటు అవుతుంది. (ఈ సంబంధం
\emph{అనుగమనం} కాదని గమనించండి.
$\mSat{M}{\Ax{T}}[w]$ అయినప్పుడల్లా
$\mSat{M}{\Ax{D}}[w]$ అవుతుందని కాదు.)
ఇలాంటి కొన్ని సంబంధాలను నమోదు చేద్దాం.

\begin{prop}\ollabel{prop:relation-facts}
  $R$ అనేది $W$ సమితిపై ద్విస్థానిక సంబంధం
  అనుకుందాం; అప్పుడు:
  \begin{enumerate}
  \item $R$ స్వావర్తనమైతే సీరియల్ కూడా.
  \item $R$ సౌష్ఠవమైతే, అది సంక్రామకమయ్యేది
    అప్పుడే, అప్పుడే అది యూక్లిడియన్ అయినప్పుడు.
  \item $R$ సౌష్ఠవం లేదా యూక్లిడియన్ అయితే,
    అది బలహీన సహగమ్యం (దానికి ``వజ్ర ధర్మం'' ఉంది).
  \item $R$ యూక్లిడియన్ అయితే అది బలహీన సంయుక్తం.
  \item $R$ ప్రమేయాత్మకమైతే అది సీరియల్.
  \end{enumerate}
\end{prop}

\begin{prob}
  \olref[nml][frd][def]{prop:relation-facts}ను నిరూపించండి.
\end{prob}

\end{document}
